%\section{Additional details on experimental setup}

Table \ref{tab:hpo_space} shows the hyper-parameter space for our experiments on the image and OpenML datasets. For our Bayesmark experiments, we use the default hyper-parameter space.

\begin{table}[t]
    \centering
    \small
    \caption{Hyper-parameters used in this study. For Bayesmark, we use the default hyper-parameter set, whose size is 134M.}
    \label{tab:hpo_space}
    \begin{tabular}{lr}
        \toprule
        \multicolumn{2}{c}{OpenML} \\
        \midrule
        \textbf{Hyper-parameter} & \textbf{Range} \\
        \midrule
        Network depth & [1, 4] \\
        Network width & [16, 32, 64, \ldots, 1024] \\
        Batch size & [4, 8, 16, \ldots, 256] \\
        Initial learning rate & $[10^{-5}, 10^{-4}, \ldots, 1.0]$ \\
        \midrule
        \multicolumn{2}{c}{MNIST} \\
        \midrule
        \textbf{Hyper-parameter} & \textbf{Range} \\
        \midrule
        Number of filters & [2, 6] \\
        Kernel size & [2, 6] \\
        Padding & Valid, same \\
        Number of conv blocks & [1, 3] \\
        \midrule
        \multicolumn{2}{c}{SVHN} \\
        \midrule
        \textbf{Hyper-parameter} & \textbf{Range} \\
        \midrule
        Number of filters & [2, 6] \\
        Kernel size & [2, 6] \\
        Padding & Valid, same \\
        Number of conv blocks & [1, 3] \\
        Dropout rate & [0.2, 0.5] \\
        Final dropout rate & [0.2, 0.5] \\
        Number of units & [32, 64, 128, \ldots, 512] \\
        \bottomrule
    \end{tabular}
\end{table}